\documentclass[tikz,border=8pt]{standalone}
\usepackage{tikz}
\usetikzlibrary{calc,arrows.meta,positioning}
\usepackage{amsmath}
\usepackage{pgfplots}
\pgfplotsset{compat=1.18}
\usepackage{ctex}

\begin{document}
\begin{tikzpicture}[
  arrstyle/.style={-{Stealth[length=9pt]}, very thick, orange!80!black},
  dotlabel/.style={font=\scriptsize, blue!80!black}
]

%% ── 左侧：n=10 个原始数据点 ──────────────────────────────────────
\node[draw, rounded corners=6pt, fill=blue!6,
      minimum width=3.2cm, minimum height=7.0cm] (databox) at (0,0) {};
\node[font=\small\bfseries, above=0.1cm of databox.north] {原始数据};
\node[font=\scriptsize, below=0.1cm of databox.south]
  {$X_1, X_2, \ldots, X_{10}$};

% 10 个数据点（固定坐标 + 手工标签）
\fill[blue!65] (-0.3, 3.0) circle (2.5pt);
\node[dotlabel, right] at (0.05, 3.0) {$x_1=2.3$};

\fill[blue!65] (-0.3, 2.3) circle (2.5pt);
\node[dotlabel, right] at (0.05, 2.3) {$x_2=1.8$};

\fill[blue!65] (-0.3, 1.6) circle (2.5pt);
\node[dotlabel, right] at (0.05, 1.6) {$x_3=3.1$};

\fill[blue!65] (-0.3, 0.9) circle (2.5pt);
\node[dotlabel, right] at (0.05, 0.9) {$x_4=2.7$};

\fill[blue!65] (-0.3, 0.2) circle (2.5pt);
\node[dotlabel, right] at (0.05, 0.2) {$x_5=1.5$};

\fill[blue!65] (-0.3,-0.5) circle (2.5pt);
\node[dotlabel, right] at (0.05,-0.5) {$x_6=2.9$};

\fill[blue!65] (-0.3,-1.2) circle (2.5pt);
\node[dotlabel, right] at (0.05,-1.2) {$x_7=2.1$};

\fill[blue!65] (-0.3,-1.9) circle (2.5pt);
\node[dotlabel, right] at (0.05,-1.9) {$x_8=3.4$};

\fill[blue!65] (-0.3,-2.6) circle (2.5pt);
\node[dotlabel, right] at (0.05,-2.6) {$x_9=1.9$};

\fill[blue!65] (-0.3,-3.3) circle (2.5pt);
\node[dotlabel, right] at (0.05,-3.3) {$x_{10}=2.6$};

% 维度标注
\node[font=\scriptsize, gray] at (-1.2, 0) {$10\,\text{维}$};

%% ── 中间箭头：充分统计量变换 ────────────────────────────────────
\node[draw, fill=orange!15, rounded corners=4pt,
      font=\small\bfseries, minimum width=2.6cm, align=center]
  (tfm) at (5.5, 0) {充分统计量\\$T = \bar{X}$};
\draw[arrstyle] (databox.east) -- (tfm.west)
  node[midway, above, font=\scriptsize]{信息浓缩};

%% ── 右侧：单个统计量 T ───────────────────────────────────────────
\node[draw, rounded corners=6pt, fill=orange!8,
      minimum width=2.6cm, minimum height=3.2cm] (statbox) at (10.0, 0) {};
\node[font=\small\bfseries, above=0.1cm of statbox.north] {统计量};
\node[font=\scriptsize, below=0.1cm of statbox.south] {$1\,\text{维}$};

% 单个数值
\fill[orange!70!black] (10.0, 0) circle (5pt);
\node[font=\normalsize, orange!60!black] at (10.0, -0.45)
  {$T = 2.43$};
\node[font=\scriptsize, gray] at (10.0, -1.0)
  {$\bar{X} = \frac{1}{n}\sum_{i=1}^n X_i$};

\draw[arrstyle] (tfm.east) -- (statbox.west);

%% ── 维度对比箭头（高亮压缩） ────────────────────────────────────
\draw[{Stealth[length=6pt]}-{Stealth[length=6pt]}, very thick, gray!60]
  ($(databox.north east)+(0.1,0.35)$) -- ($(statbox.north west)+(-0.1,0.35)$)
  node[midway, above, font=\scriptsize, gray!80]{$10\,\text{维} \longrightarrow 1\,\text{维}$};

%% ── 下方：充分性条件公式 ─────────────────────────────────────────
\node[draw, fill=green!8, rounded corners=6pt,
      font=\small, align=center, minimum width=14cm]
  (fmbox) at (5.0, -5.8) {
    \textbf{充分性条件（因子分解定理）：}\\[4pt]
    $P(X_1,\ldots,X_n \mid T=t,\,\theta)
      = P(X_1,\ldots,X_n \mid T=t)$
    \quad（与参数 $\theta$ 无关）\\[4pt]
    \small 即：给定充分统计量 $T$，样本的条件分布不含关于 $\theta$ 的额外信息
  };

%% ── 整体标题 ────────────────────────────────────────────────────
\node[font=\large\bfseries]
  at ($(databox.north)!0.5!(statbox.north)+(0,1.3cm)$)
  {充分统计量：数据信息的无损压缩};

\end{tikzpicture}
\end{document}
